\chapter{Estratégia de Implementação e Evolução do Âmbito}

\section{Da Visão Inicial à V1}

A visão inicial do NatureProtector era mais ampla do que a entrega final da V1. O projeto partiu da necessidade de apoiar a prevenção de incêndios através de observação ambiental, avaliação de risco, visualização operacional e mecanismos de apoio à decisão. Numa versão mais madura, o sistema poderia integrar sensores reais, dados territoriais mais completos, modelos calibrados, painéis operacionais, workflows de alerta e comparação com referências reconhecidas de perigo de incêndio, como o FWI, produtos nacionais de perigo de incêndio e plataformas europeias de monitorização~\citep{nrcan_fwi_system,ipma_fwi,ipma_pir_rcm,ec_jrc_effis}.

Contudo, tentar entregar uma plataforma completa de prevenção de incêndios num primeiro ciclo tornaria o projeto demasiado extenso e difícil de validar. Uma solução operacional completa exigiria dados reais, calibração científica, validação com ocorrências históricas, implantação em produção, segurança, gestão de utilizadores e integração institucional. Estes elementos são relevantes, mas excedem aquilo que podia ser fechado com rigor na V1.

Por esse motivo, a estratégia de implementação foi delimitada para uma base técnica do módulo de prevenção. A V1 concentra-se na cadeia essencial que sustenta a visão futura: simulação de observações, transporte de eventos, processamento durável, avaliação candidata de risco, projeções operacionais, exposição por API, observabilidade e recolha de evidência. A utilização de transporte assíncrono e processamento durável é coerente com práticas comuns em sistemas orientados a mensagens, nos quais a entrega, o consumo, a confirmação e o processamento efetivo devem ser tratados como etapas distintas~\citep{rabbitmq_amqp_model,rabbitmq_consumer_acknowledgements,richardson_idempotent_consumer}.

Esta delimitação torna a entrega mais defensável. A V1 não é apresentada como modelo científico final, índice oficial ou solução pronta para produção. É apresentada como uma fundação técnica executável, rastreável e auditável, sobre a qual podem ser construídas fases futuras de calibração, validação e integração operacional.

\section{Âmbito Consolidado}

A consolidação do âmbito foi uma das decisões centrais do projeto. A entrega final foi concentrada no módulo de prevenção e na baseline local necessária para o demonstrar.

O âmbito consolidado inclui leituras simuladas de sensores, publicação e consumo de eventos por RabbitMQ, persistência durável em PostgreSQL, validação e elegibilidade das leituras, tratamento de rejeições e quarentenas, geração de avaliações candidatas de risco, projeções operacionais por área e por célula, política interna de alertas, exposição através da Backoffice API, observabilidade com InfluxDB e Grafana, testes automatizados e recolha de evidência em tempo de execução.

Este âmbito permite demonstrar uma cadeia técnica de ponta a ponta, desde a produção de leituras simuladas até à exposição de estado operacional. Ao mesmo tempo, evita tratar como concluídos aspetos que exigem trabalho científico ou operacional adicional.

O princípio orientador foi o seguinte:

\begin{quote}
A V1 deve fechar uma base técnica e metodológica antes de afirmar maturidade científica ou operacional.
\end{quote}

Consequentemente, ficaram fora da entrega atual a calibração científica, a validação contra ocorrências reais, a equivalência com sistemas oficiais, a integração com sensores físicos reais, a implantação em produção e workflows operacionais completos de alerta. Estes elementos não são falhas da V1; são limites deliberados do seu âmbito.

\section{Fases de Implementação}

A implementação evoluiu através de fases práticas que clarificaram a arquitetura, os contratos, a pipeline e a evidência. A Tabela~\ref{tab:implementation-phases} resume essa evolução.

\begin{table}[h]
\centering
\caption{Fases principais de implementação da V1.}
\label{tab:implementation-phases}
\begin{tabular}{p{0.27\textwidth} p{0.56\textwidth}}
\hline
\textbf{Fase} & \textbf{Resultado principal} \\
\hline
Infraestrutura local & Integração de RabbitMQ, PostgreSQL, InfluxDB e Grafana. \\
Simulação e controlo & Configuração da área piloto, sensores, grelha, cenários e execuções simuladas. \\
Contratos e eventos & Definição da fronteira baseada em \texttt{EventEnvelope<TPayload>} e \texttt{SensorReadingProduced}. \\
Pipeline durável & Persistência de eventos em inbox, registo de tentativas, rejeições e quarentenas. \\
Elegibilidade e risco & Separação entre leitura normalizada, qualidade, elegibilidade, \texttt{RiskInput} e \texttt{RiskAssessment}. \\
Projeções, API e alertas & Atualização de estados operacionais, exposição pela API e política interna com \texttt{None}, \texttt{Warning} e \texttt{Alarm}. \\
Validação C7 & Recolha de evidência em tempo de execução, comparação API/base de dados e execução de testes automatizados. \\
\hline
\end{tabular}
\end{table}

Estas fases mostram que a implementação não se limitou à criação de uma fórmula de risco. O trabalho principal consistiu em criar uma cadeia técnica rastreável, com fronteiras claras entre transporte, validação, elegibilidade, avaliação, projeção, exposição e validação.
